%==============================================================================
%  CHAPTER 1 — GENERAL CONTEXT OF THE PROJECT
%==============================================================================
\chapter{General Context of the Project}
\label{chap:context}

\section*{Introduction}
This chapter sets the stage for the project. It presents the host organization,
the framework of the end-of-studies project, a study of the existing situation
and its shortcomings, the proposed solution, and finally the method and
work environment adopted to carry out the work.

% ---------------------------------------------------------------------------
\section{Host Organization}
\label{sec:host-org}

% >>> USER: replace this section with INSOMEA's real presentation. Keep the
% >>> sub-structure (presentation / activities / org chart). Add the logo and
% >>> an org-chart figure into images/ and reference it.

\textbf{INSOMEA} (INSOMEA Computer Solutions) is a Microsoft Cloud Gold Partner
and one of the leading cloud IT solution providers in the Middle East and Africa
(MEA) region. Founded in 2016, it fulfills the IT requirements of local, regional,
and international organizations and specializes in Microsoft~365, Microsoft Azure,
cybersecurity, and professional training. The company operates from six offices
across Tunisia, Bahrain, France, Algeria, Morocco, and the Ivory Coast, with a
team of more than fifty employees serving over five hundred customers, and more
than eight years of expertise. A Microsoft Cloud Partner since 2017, INSOMEA has
been named Microsoft \emph{Partner of the Year} on five occasions, in Bahrain
(2020) and in Tunisia (2021 and 2022) among them, and holds more than fifteen
Gold competencies together with eight advanced specializations.

\subsection{Areas of Activity}
INSOMEA delivers smart, cost-effective IT solutions built on Microsoft
technologies, organized around three main pillars:
\begin{itemize}[leftmargin=1.4em]
  \item \textbf{Modern Workplace}, modernizing collaboration and productivity
        with Microsoft~365 and Microsoft Teams, including adoption and
        change-management services;
  \item \textbf{Security}, securing client environments through identity and
        access management, information protection and governance, and threat
        protection;
  \item \textbf{Cloud \& Infrastructure}, Microsoft Azure cloud platform,
        datacenter, DevOps, and application development / integration.
\end{itemize}
In addition, the \textbf{INSOMEA Training Center} offers a broad catalog of
Microsoft certification courses delivered by Microsoft Certified Trainers (MCTs).
The company's engagement methodology rests on four core values: end-to-end change
management, a fully certified technical team, solution adoption, and continuous
support.

% ---------------------------------------------------------------------------
\section{Project Framework}
\label{sec:framework}
This work was conducted as an end-of-studies project (Projet de Fin d'\'Etudes),
the final milestone of the engineering cycle, over a period of six months. Its
objective is to design and build a complete \gls{mlops} platform that automates
the machine-learning lifecycle, from code upload through to a live, queryable
prediction service.

% ---------------------------------------------------------------------------
\section{Study of the Existing Situation}
\label{sec:existing}
Before designing the solution, it is useful to characterize how machine-learning
models are typically taken from experiment to production in teams that lack a
dedicated platform. The prevailing workflow is largely manual:

\begin{enumerate}[leftmargin=1.6em]
  \item A data scientist trains a model locally, inside a Jupyter notebook, on a
        personal machine.
  \item Results (accuracy, loss, hyper-parameters) are recorded informally ---
        in the notebook output, a spreadsheet, or not at all.
  \item The trained model is serialized to a file and shared by email or a shared
        drive.
  \item An engineer wraps the model in a web service by hand and deploys it
        through custom scripts onto a server.
  \item Monitoring, if any, is limited to checking that the service responds.
\end{enumerate}

% ---------------------------------------------------------------------------
\section{Critique of the Existing Situation}
\label{sec:critique}
This manual workflow suffers from well-known weaknesses:

\begin{itemize}[leftmargin=1.4em]
  \item \textbf{No reproducibility.} Without a record of the exact code, data,
        and parameters of a run, results cannot be reliably reproduced.
  \item \textbf{Lost experiments.} Informal tracking means that comparisons
        between models are unreliable and historical runs are quickly forgotten.
  \item \textbf{``Works on my machine''.} Differences between the developer's
        environment and the production server cause failures that are hard to
        diagnose.
  \item \textbf{Slow, error-prone deployment.} Hand-written deployment scripts
        are brittle, must be repeated for every model, and are a frequent source
        of incidents.
  \item \textbf{No scalability.} A single hand-managed server cannot elastically
        scale to demand or isolate models from one another.
\end{itemize}

These limitations motivate a platform that automates and standardizes the entire
lifecycle.

% ---------------------------------------------------------------------------
\section{Proposed Solution}
\label{sec:solution}
The proposed solution is a self-hosted \gls{mlops} platform that lets a single
user, through a web browser:

\begin{enumerate}[leftmargin=1.6em]
  \item sign up and sign in (token-based authentication);
  \item create a \textbf{project} as an organizing unit, each backed by its own
        experiment-tracking namespace;
  \item upload code (\texttt{.py}, \texttt{.ipynb}, or \texttt{.zip}) and a
        dataset through the browser;
  \item trigger a \textbf{training pipeline} that runs on a \gls{k8s} cluster;
  \item automatically capture metrics, parameters, and artifacts from each run;
  \item browse, \textbf{register}, and \textbf{promote} model versions
        (Staging / Production / Archived);
  \item \textbf{deploy} a promoted model as a live \gls{http} prediction service;
  \item send sample inputs and receive predictions back;
  \item monitor live cluster and per-pod metrics from a dashboard;
  \item (as an administrator) manage user accounts and roles and oversee the
        whole platform.
\end{enumerate}

The platform runs on a real Kubernetes cluster and ships every change to a public
URL through an automated \gls{cicd} pipeline.

% ---------------------------------------------------------------------------
\section{Adopted Method: Scrum}
\label{sec:scrum}
Given that the platform is built incrementally from a set of independent but
connected capabilities, an \textbf{Agile} approach was chosen over a classic
sequential (waterfall) process. Agile favors short iterations, frequent working
increments, and the ability to re-prioritize, a natural fit for a project
whose features (auth, pipelines, registry, deployment, monitoring) can be
delivered and demonstrated one layer at a time.

Among Agile methods, \textbf{Scrum}~\cite{scrum_guide} was adopted. Scrum is an
iterative and incremental Agile method that organizes work into fixed-length
iterations called \emph{sprints}, each producing a potentially shippable
increment. Figure~\ref{fig:scrum} illustrates the overall Scrum process.

\begin{figure}[H]
  \centering
  \includegraphics[width=0.9\textwidth,keepaspectratio]{scrum}
  \caption{The Scrum process: from product backlog to shippable increment}
  \label{fig:scrum}
\end{figure}

\subsection{Scrum Roles}
\begin{itemize}[leftmargin=1.4em]
  \item \textbf{\gls{po}}: Mr.\ Amine Gonji (company supervisor, INSOMEA),
        responsible for prioritizing the product backlog.
  \item \textbf{\gls{sm}}: Mr.\ Ben Mardes Achref (academic supervisor),
        responsible for facilitating the process and removing impediments.
  \item \textbf{Development Team}: the student (author), responsible for analysis,
        design, implementation, and testing.
\end{itemize}

\subsection{Scrum Artifacts}
\begin{itemize}[leftmargin=1.4em]
  \item \textbf{Product Backlog}: the ordered list of all desired features,
        expressed as user stories (see Chapter~\ref{chap:requirements}).
  \item \textbf{Sprint Backlog}: the subset of stories selected for a given
        sprint, broken into tasks.
  \item \textbf{Increment}: the working, demonstrable software produced at the end
        of each sprint.
\end{itemize}

\subsection{Scrum Events}
The project was organized into \textbf{four sprints}. Each sprint began with
\emph{sprint planning}, proceeded with regular progress reviews, and ended with a
\emph{sprint review} (demonstration of the increment) and a \emph{retrospective}
(lessons learned). The detailed release plan is given in
Chapter~\ref{chap:requirements}.

% ---------------------------------------------------------------------------
\section{Work Environment}
\label{sec:environment}

\subsection{Hardware Environment}
% >>> USER: adjust the development-machine row to your real specs.
\begin{table}[H]\centering
\renewcommand{\arraystretch}{1.3}
\caption{Hardware environment}\label{tab:hardware}
\begin{tabularx}{\textwidth}{|>{\raggedright\arraybackslash}p{4.2cm}|X|}
\hline
\textbf{Component} & \textbf{Specification}\\
\hline
Development machine & Intel Core i7, 16\,GB RAM, Windows 11\\
\hline
Cloud cluster       & Azure Kubernetes Service (\gls{aks})\\
\hline
\end{tabularx}
\end{table}

\subsection{Software Environment and Tools}
\begin{table}[H]\centering
\renewcommand{\arraystretch}{1.3}
\caption{Software tools and technologies}\label{tab:software}
\begin{tabularx}{\textwidth}{|>{\raggedright\arraybackslash}p{4.2cm}|X|}
\hline
\textbf{Category} & \textbf{Tools / Technologies}\\
\hline
Backend             & FastAPI (Python~3.11), SQLAlchemy~2 (async), Pydantic~v2\\
\hline
Frontend            & Angular~19 (standalone components, signals), Tailwind~v4\\
\hline
Database            & PostgreSQL~16\\
\hline
Object storage      & MinIO (\gls{s3}-compatible)\\
\hline
Experiment tracking & MLflow\\
\hline
Orchestration       & Kubernetes (\gls{aks} / KinD), \gls{kfp}, Argo Workflows\\
\hline
Model serving       & KServe\\
\hline
Networking          & Cloudflare Tunnel (\gls{tls} termination)\\
\hline
CI/CD               & GitHub Actions, Azure Container Registry (\gls{acr})\\
\hline
Version control     & Git / GitHub\\
\hline
Modeling            & \gls{uml} (draw.io / PlantUML)\\
\hline
IDE                 & Visual Studio Code\\
\hline
\end{tabularx}
\end{table}

\section*{Conclusion}
This chapter introduced the host organization, the project framework, and the
shortcomings of the manual machine-learning workflow that motivate this work. It
presented the proposed \gls{mlops} platform at a high level, and the Scrum
method and tooling adopted. The next chapter establishes the technical
background and justifies the technology choices in depth.
